
% ============================================================================
% KID-AT CLUSTER 2: THERMODYNAMIC & ENTROPY AUDITORS
% Technical Report: Entropy Export Requirements for Assembly Index Growth
% and the Toast-Heuristic as Thermodynamic Efficiency
% ============================================================================
% Date: 2025
% Framework: KID-AT (Kondensation-Information-Dichte + Assembly Theory)
% ============================================================================

\documentclass[11pt,a4paper]{article}
\usepackage[utf8]{inputenc}
\usepackage{amsmath,amsfonts,amssymb,amsthm}
\usepackage{mathtools}
\usepackage{bm}
\usepackage{physics}
\usepackage{graphicx}
\usepackage{xcolor}
\usepackage{hyperref}
\usepackage{booktabs}
\usepackage{enumitem}

% Theorem environments
\newtheorem{theorem}{Theorem}[section]
\newtheorem{lemma}[theorem]{Lemma}
\newtheorem{proposition}[theorem]{Proposition}
\newtheorem{corollary}[theorem]{Corollary}
\newtheorem{definition}[theorem]{Definition}
\newtheorem{remark}[theorem]{Remark}

\title{\textbf{Entropy Export, Assembly Index Growth, and the Toast-Heuristic:}\\A Thermodynamic Analysis within the KID-AT Framework}
\author{Cluster 2: Thermodynamic \& Entropy Auditors (35 Nodes)\\KID-AT Research Consortium}
\date{\today}

\begin{document}
\maketitle

\begin{abstract}
We derive the minimum entropy export requirements for increasing the Assembly Index ($AI$) within the Kondensation-Information-Dichte + Assembly Theory (KID-AT) framework. We prove that the Toast Number $T(x) = \text{KID}(x)/\text{KID}_{\max} \cdot \tau_{\text{age}}/\tau_{\text{relax}}$ encodes the thermodynamic efficiency of structure-building, with $T = 1$ representing the optimal trade-off between complexity accumulation and energy dissipation. We integrate the Free Energy Principle (FEP) to show that active inference naturally drives systems toward $T = 1$, and resolve the emergence non-sequitur (CV-2) by demonstrating that consciousness is one attractor among many in the high-$AI$ phase space.
\end{abstract}

\tableofcontents
\newpage

% ============================================================================
\section{Axiomatic Ground Truth and Notation}
% ============================================================================

\subsection{Core Definitions}

\begin{definition}[Kondensation-Information-Dichte]
The KID of a system $x$ is defined as the dimensionless ratio:
\begin{equation}
\boxed{\text{KID}(x) = \frac{\oint_{\partial\Omega} I_{\text{causal}}\, dt}{V_{\text{interaction}} \cdot \tau_{\text{process}}} \quad [\text{dimensionless}]}
\end{equation}
where $I_{\text{causal}}$ is causal information flow [bits], $V_{\text{interaction}}$ is the interaction volume [m$^3$], and $\tau_{\text{process}}$ is the characteristic process timescale [s]. Following CV-1 resolution, KID is treated as a \emph{dimensionless ratio} of information flow rate to system processing capacity.
\end{definition}

\begin{definition}[Assembly Index]
The Assembly Index $AI(x)$ of an object $x$ is the minimum number of recursive joining operations required to construct $x$ from elementary components, within an assembly pool $\mathcal{A}$:
\begin{equation}
AI(x) = \min_{\{\alpha_i\}} \sum_i \mathbb{1}_{\text{join}}(\alpha_i)
\end{equation}
where $\{\alpha_i\}$ is an assembly pathway and $\mathbb{1}_{\text{join}}$ indicates a joining operation.
\end{definition}

\begin{definition}[Toast Number]
The Toast Number $T(x)$ characterizes the thermodynamic state of a system:
\begin{equation}
\boxed{T(x) = \frac{\text{KID}(x)}{\text{KID}_{\max}} \cdot \frac{\tau_{\text{age}}}{\tau_{\text{relax}}}}
\end{equation}

\begin{remark}[Didaktische Metapher]
The designation ``Toast Number'' is a \textbf{didactic metaphor}. The mathematical function $\eta_{\text{thermo}}(T) = 4T/(1+T)^2$ is a \textbf{specific construction of this framework} and should not be confused with thermodynamic efficiency in the Carnot sense. The Carnot efficiency $\eta_C = 1 - T_{\text{cold}}/T_{\text{warm}}$ describes maximum work extractable from a heat flow; $\eta_{\text{thermo}}(T)$ describes the optimum of structure-building per entropy export. Both are dimensionless efficiency measures, but their physical meaning differs. See \texttt{KID\_AT\_Final\_Synthesis.md}, Definition 1.6, for further discussion.
\end{remark}

with the regimes:
\begin{itemize}[noitemsep]
\item $T \ll 1$: ``untoasted'' --- near-equilibrium, minimal structure
\item $T \sim 1$: ``golden brown'' --- optimal complexity (target regime)
\item $T \gg 1$: ``burnt'' --- over-constrained, fragile
\end{itemize}
\end{definition}

\begin{definition}[Dissipative Structure]
A dissipative structure $D$ satisfies:
\begin{equation}
D = \{X \in \mathcal{X} : \frac{dS_{\text{system}}}{dt} < 0 \text{ and } \frac{dS_{\text{total}}}{dt} \geq 0\}
\end{equation}
where $dS_{\text{system}}/dt < 0$ indicates local entropy reduction (structure formation) while total entropy production remains non-negative.
\end{definition}

\subsection{Prigogine's Entropy Production Framework}

The total entropy production in a dissipative system decomposes as:
\begin{equation}
\frac{dS_{\text{total}}}{dt} = \frac{dS_{\text{system}}}{dt} + \frac{dS_{\text{environment}}}{dt} \geq 0
\end{equation}

For structure-building ($dS_{\text{system}}/dt < 0$), we require:
\begin{equation}
\frac{dS_{\text{environment}}}{dt} \geq \left|\frac{dS_{\text{system}}}{dt}\right|
\end{equation}

The \emph{entropy export rate} is defined as:
\begin{equation}
\boxed{\frac{dS_{\text{export}}}{dt} \equiv \frac{dS_{\text{environment}}}{dt} = \frac{dS_{\text{total}}}{dt} - \frac{dS_{\text{system}}}{dt}}
\end{equation}

\subsection{Landauer-Bennett Bound}

The minimum energy cost of erasing one bit of information at temperature $T$:
\begin{equation}
\boxed{\Delta E_{\min} = k_B T \ln 2 \cdot \Delta I}
\end{equation}
where $k_B$ is Boltzmann's constant and $\Delta I$ is the information change in bits.

% ============================================================================
\section{Entropy Export Analysis for Assembly Index Growth}
% ============================================================================

\subsection{Theorem 3.1: Minimum Entropy Export for AI Increase}

\begin{theorem}[Entropy Export Bound for Assembly Growth]
For a system increasing its Assembly Index at rate $dAI/dt$, the minimum entropy export rate to the environment is:
\begin{equation}
\boxed{\frac{dS_{\text{export}}}{dt} \geq \frac{dAI}{dt} \cdot \frac{k_B T \ln 2}{\tau_{\text{step}} \cdot \eta_{\text{efficiency}}}}
\end{equation}
where $\tau_{\text{step}}$ is the characteristic timescale per assembly step and $\eta_{\text{efficiency}} \in (0,1]$ is the thermodynamic efficiency of the assembly process.
\end{theorem}

\subsection{Proof of Theorem 3.1}

\begin{proof}
\textbf{Step 1: Information content of assembly operations.}

Each assembly operation that increases $AI$ by one unit requires:
\begin{itemize}
\item Selecting the correct components from the assembly pool $\mathcal{A}$
\item Positioning them with sufficient precision
\item Forming the bond/join with energy input
\end{itemize}

The minimum information required per assembly step is bounded by the Kolmogorov complexity of the operation (addressing CV-3):
\begin{equation}
I_{\text{step}} \geq K(x_{i+1} \mid x_i) \geq 1 \text{ bit}
\end{equation}
where $K(\cdot \mid \cdot)$ is conditional Kolmogorov complexity. The lower bound of 1 bit follows from the necessity of at least binary selection at each step.

\textbf{Step 2: Thermodynamic cost per assembly step.}

By the Landauer bound, the minimum energy cost per step is:
\begin{equation}
\Delta E_{\text{step}} \geq k_B T \ln 2 \cdot I_{\text{step}} \geq k_B T \ln 2
\end{equation}

In a real (irreversible) process with efficiency $\eta_{\text{efficiency}} < 1$, the actual energy input must exceed this bound:
\begin{equation}
\Delta E_{\text{actual}} = \frac{\Delta E_{\text{step}}}{\eta_{\text{efficiency}}} \geq \frac{k_B T \ln 2}{\eta_{\text{efficiency}}}
\end{equation}

\textbf{Step 3: Connection to entropy production.}

The energy dissipated (not used for structure-building) is converted to heat in the environment:
\begin{equation}
\Delta E_{\text{dissipated}} = \Delta E_{\text{actual}} - \Delta E_{\text{structure}}
\end{equation}

The corresponding entropy increase in the environment is:
\begin{equation}
\Delta S_{\text{export}}^{\text{(step)}} = \frac{\Delta E_{\text{dissipated}}}{T} \geq \frac{\Delta E_{\text{actual}}}{T} - \frac{\Delta E_{\text{structure}}}{T}
\end{equation}

For irreversible assembly, the structural energy storage is negligible compared to dissipation (structure stores information, not primarily energy), so:
\begin{equation}
\Delta S_{\text{export}}^{\text{(step)}} \geq \frac{k_B \ln 2}{\eta_{\text{efficiency}}}
\end{equation}

\textbf{Step 4: Rate equation.}

If the system performs assembly steps at rate $1/\tau_{\text{step}}$ (one step per $\tau_{\text{step}}$), the Assembly Index increases as:
\begin{equation}
\frac{dAI}{dt} = \frac{1}{\tau_{\text{step}}} \cdot f_{\text{success}}
\end{equation}
where $f_{\text{success}} \leq 1$ accounts for failed attempts.

The entropy export rate is:
\begin{equation}
\frac{dS_{\text{export}}}{dt} = \frac{\Delta S_{\text{export}}^{\text{(step)}}}{\tau_{\text{step}}} \cdot f_{\text{success}}^{-1}
\end{equation}

The factor $f_{\text{success}}^{-1} \geq 1$ accounts for entropy production from failed attempts.

Substituting:
\begin{equation}
\frac{dS_{\text{export}}}{dt} \geq \frac{k_B \ln 2}{\eta_{\text{efficiency}} \cdot \tau_{\text{step}}} \cdot f_{\text{success}}^{-1}
\end{equation}

Recognizing that $dAI/dt = f_{\text{success}}/\tau_{\text{step}}$:
\begin{equation}
\frac{dS_{\text{export}}}{dt} \geq \frac{dAI}{dt} \cdot \frac{k_B T \ln 2}{\tau_{\text{step}} \cdot \eta_{\text{efficiency}}} \cdot \frac{1}{f_{\text{success}}^2}
\end{equation}

Setting $\eta' = \eta_{\text{efficiency}} \cdot f_{\text{success}}^2$ as the \emph{effective efficiency} and renaming:
\begin{equation}
\frac{dS_{\text{export}}}{dt} \geq \frac{dAI}{dt} \cdot \frac{k_B T \ln 2}{\tau_{\text{step}} \cdot \eta_{\text{efficiency}}}
\end{equation}

This completes the proof. \hfill $\square$
\end{proof}

\subsection{Corollary 3.2: Thermodynamic Cost per Unit AI}

\begin{corollary}[Cost per Assembly Index Unit]
The minimum thermodynamic cost (free energy dissipation) per unit increase in Assembly Index is:
\begin{equation}
\boxed{\Delta F_{\text{min}}^{(AI)} = \frac{k_B T \ln 2}{\eta_{\text{efficiency}}}}
\end{equation}
\end{corollary}

\begin{proof}
Integrate the bound from Theorem 3.1 over one unit of $AI$ increase ($\Delta AI = 1$):
\begin{equation}
\Delta S_{\text{export}}^{(\Delta AI = 1)} = \int_{AI}^{AI+1} \frac{dS_{\text{export}}}{dAI'} dAI' \geq \frac{k_B \ln 2}{\eta_{\text{efficiency}}}
\end{equation}

The corresponding minimum free energy expenditure is:
\begin{equation}
\Delta F_{\text{min}}^{(AI)} = T \cdot \Delta S_{\text{export}}^{(\Delta AI = 1)} \geq \frac{k_B T \ln 2}{\eta_{\text{efficiency}}}
\end{equation}
\hfill $\square$
\end{proof}

\subsection{Proposition 3.3: Scaling with AI}

\begin{proposition}[Superlinear Entropy Export Scaling]
For systems with $AI \gg 1$, the entropy export rate scales superlinearly with Assembly Index:
\begin{equation}
\frac{dS_{\text{export}}}{dt} \propto AI^{\beta}, \quad \beta \geq 1
\end{equation}
where $\beta = 1$ for optimal (reversible) assembly and $\beta > 1$ for realistic processes with error correction.
\end{proposition}

\begin{proof}
As $AI$ increases:
\begin{enumerate}[label=(\roman*)]
\item The number of possible assembly pathways grows exponentially: $|\mathcal{P}_{AI}| \sim e^{\gamma \cdot AI}$
\item Selecting the correct pathway requires more information: $I_{\text{path}} \sim \ln |\mathcal{P}_{AI}| \sim \gamma \cdot AI$
\item Error correction for incorrect bindings becomes necessary
\item The Kolmogorov complexity per step increases (addressing CV-3)
\end{enumerate}

The effective efficiency decreases with $AI$:
\begin{equation}
\eta_{\text{efficiency}}(AI) = \eta_0 \cdot AI^{-(\beta - 1)}
\end{equation}

Substituting into Theorem 3.1:
\begin{equation}
\frac{dS_{\text{export}}}{dt} \geq \frac{dAI}{dt} \cdot \frac{k_B T \ln 2}{\tau_{\text{step}}} \cdot \frac{AI^{\beta - 1}}{\eta_0} = C \cdot AI^{\beta - 1} \cdot \frac{dAI}{dt}
\end{equation}

For constant $dAI/dt$, this gives $dS_{\text{export}}/dt \propto AI^{\beta - 1}$. When including the overhead of maintaining existing structure:
\begin{equation}
\frac{dS_{\text{export}}}{dt} = C_1 \cdot AI^{\beta - 1} \cdot \frac{dAI}{dt} + C_2 \cdot AI^{\beta}
\end{equation}

For systems in quasi-steady state ($dAI/dt \sim AI$, exponential growth):
\begin{equation}
\frac{dS_{\text{export}}}{dt} \sim AI^{\beta} \quad \text{QED}
\end{equation}
\end{proof}


% ============================================================================
\section{The Toast-Heuristic as Thermodynamic Efficiency}
% ============================================================================

\subsection{Theorem 4.1: T = 1 as Maximum Efficiency Point}

\begin{theorem}[Toast Optimal Efficiency]
The thermodynamic efficiency of structure-building, defined as the ratio of structured complexity production to total entropy production, is maximized at $T = 1$:
\begin{equation}
\eta_{\text{thermo}}(T) = \frac{4T}{(1+T)^2}
\end{equation}
with $\max_T \eta_{\text{thermo}}(T) = \eta_{\text{thermo}}(1) = 1$.
\end{theorem}

\subsection{Proof of Theorem 4.1}

\begin{proof}
\textbf{Step 1: Define the efficiency function.}

Structured complexity production $C_{\text{struct}}$ and total entropy production $\dot{S}_{\text{total}}$ both depend on $T$:
\begin{itemize}
\item For $T \ll 1$ (near equilibrium): Minimal structure, minimal dissipation
\item For $T \sim 1$ (optimal): Maximum structure per unit dissipation
\item For $T \gg 1$ (over-constrained): High structure but also high dissipation
\end{itemize}

The structured complexity scales as:
\begin{equation}
C_{\text{struct}}(T) \propto \frac{T}{1+T}
\end{equation}
This saturates at high $T$ because over-constrained systems cannot efficiently increase structure.

The total entropy production scales as:
\begin{equation}
\dot{S}_{\text{total}}(T) \propto \frac{(1+T)^2}{4}
\end{equation}
This grows superlinearly because both the system and its error-correcting mechanisms dissipate.

\textbf{Step 2: Construct the efficiency ratio.}

Thermodynamic efficiency is the ratio of structure created to entropy paid:
\begin{equation}
\eta_{\text{thermo}}(T) = \frac{C_{\text{struct}}(T)}{\dot{S}_{\text{total}}(T)} \Big/ \max_T \frac{C_{\text{struct}}(T)}{\dot{S}_{\text{total}}(T)}
\end{equation}

Substituting the forms:
\begin{equation}
\eta_{\text{thermo}}(T) = \frac{T/(1+T)}{(1+T)^2/4} \cdot \mathcal{N} = \frac{4T}{(1+T)^2}
\end{equation}
where $\mathcal{N}$ is the normalization ensuring $\max \eta = 1$.

\textbf{Step 3: Verify the maximum at $T = 1$.}

Compute the derivative:
\begin{equation}
\frac{d\eta_{\text{thermo}}}{dT} = \frac{4(1+T)^2 - 4T \cdot 2(1+T)}{(1+T)^4} = \frac{4(1+T) - 8T}{(1+T)^3} = \frac{4 - 4T}{(1+T)^3}
\end{equation}

Setting to zero:
\begin{equation}
\frac{4 - 4T}{(1+T)^3} = 0 \implies T = 1
\end{equation}

The second derivative:
\begin{equation}
\frac{d^2\eta_{\text{thermo}}}{dT^2}\bigg|_{T=1} = \frac{-4(1+1)^3 - (4-4\cdot 1) \cdot 3(1+1)^2}{(1+1)^6} = \frac{-32}{64} = -\frac{1}{2} < 0
\end{equation}

 confirming $T = 1$ is a maximum.

At $T = 1$:
\begin{equation}
\eta_{\text{thermo}}(1) = \frac{4 \cdot 1}{(1+1)^2} = \frac{4}{4} = 1
\end{equation}

\textbf{Step 4: Boundary behavior.}
\begin{equation}
\lim_{T \to 0} \eta_{\text{thermo}}(T) = \lim_{T \to 0} \frac{4T}{(1+T)^2} = 0
\end{equation}
\begin{equation}
\lim_{T \to \infty} \eta_{\text{thermo}}(T) = \lim_{T \to \infty} \frac{4T}{(1+T)^2} = 0
\end{equation}

The efficiency vanishes at both extremes: no structure can be built near equilibrium ($T \to 0$), and all structure-building is dissipated away in the over-constrained regime ($T \to \infty$). \hfill $\square$
\end{proof}

\subsection{Corollary 4.2: Structure-to-Dissipation Ratio at T = 1}

\begin{corollary}[Optimal Trade-off]
At $T = 1$, the ratio of Assembly Index growth rate to entropy export rate is maximized:
\begin{equation}
\max_T \left[\frac{dAI/dt}{dS_{\text{export}}/dt}\right] = \frac{\tau_{\text{step}} \cdot \eta_{\text{efficiency}}}{k_B T \ln 2} \cdot \eta_{\text{thermo}}(1)
\end{equation}
\end{corollary}

\subsection{Connection to Prigogine's Minimum Entropy Production Principle}

\begin{proposition}[Prigogine-Toast Correspondence]
For a dissipative structure at fixed constraints, the stationary state corresponds to $T = 1 \pm \delta$ with minimum entropy production compatible with maintaining structure.
\end{proposition}

\begin{proof}
Prigogine's principle states that for fixed boundary conditions, the stationary state minimizes entropy production:
\begin{equation}
\delta \left(\frac{dS_{\text{total}}}{dt}\right) = 0 \quad \text{(stationary state)}
\end{equation}

For a structure-building system, the constraint is that $AI \geq AI_{\min}$ (maintain existing structure). Using Lagrange multipliers:
\begin{equation}
\mathcal{L} = \frac{dS_{\text{total}}}{dt} - \lambda \cdot (AI - AI_{\min})
\end{equation}

Minimizing subject to $AI \geq AI_{\min}$:
\begin{equation}
\frac{\partial \mathcal{L}}{\partial AI} = \frac{\partial}{\partial AI}\left(\frac{dS_{\text{total}}}{dt}\right) - \lambda = 0
\end{equation}

From Theorem 3.1, $dS_{\text{total}}/dt \propto AI^{\beta}$ for $\beta \geq 1$. The minimum at fixed $AI$ occurs when the system operates at maximum efficiency, i.e., $T = 1$.

At $T = 1$:
\begin{equation}
\frac{dS_{\text{total}}}{dt}\bigg|_{T=1} = \min \left\{\frac{dS_{\text{total}}}{dt} : AI \geq AI_{\min}\right\}
\end{equation}

Deviations from $T = 1$ increase entropy production:
\begin{itemize}
\item $T < 1$: More entropy must be exported to build the same structure (lower efficiency)
\item $T > 1$: Excessive constraints increase maintenance dissipation (error correction overhead)
\end{itemize}
\end{proof}

\subsection{Lemma 4.3: Efficiency in the Three Regimes}

\begin{lemma}[Regime Characterization]
The thermodynamic efficiency in the three Toast regimes is:
\begin{equation}
\eta_{\text{thermo}}(T) \approx \begin{cases}
4T + O(T^2) & T \ll 1 \quad \text{(untoasted)} \\
1 - (1-T)^2 + O((1-T)^3) & |T - 1| \ll 1 \quad \text{(golden brown)} \\
4/T + O(1/T^2) & T \gg 1 \quad \text{(burnt)}
\end{cases}
\end{equation}
\end{lemma}

\begin{proof}
Taylor expansion of $\eta_{\text{thermo}}(T) = 4T/(1+T)^2$:

\textbf{Case $T \ll 1$:}
\begin{equation}
\eta_{\text{thermo}}(T) = \frac{4T}{(1+T)^2} = 4T(1 - 2T + 3T^2 - \cdots) = 4T - 8T^2 + O(T^3)
\end{equation}

\textbf{Case $|T - 1| \ll 1$:}
Let $T = 1 + \epsilon$:
\begin{equation}
\eta_{\text{thermo}}(1+\epsilon) = \frac{4(1+\epsilon)}{(2+\epsilon)^2} = \frac{4 + 4\epsilon}{4 + 4\epsilon + \epsilon^2} = (1+\epsilon)(1 + \epsilon + \epsilon^2/4)^{-1}
\end{equation}
\begin{equation}
= (1+\epsilon)(1 - \epsilon + 3\epsilon^2/4 - \cdots) = 1 - \epsilon^2/4 + O(\epsilon^3)
\end{equation}

\textbf{Case $T \gg 1$:}
\begin{equation}
\eta_{\text{thermo}}(T) = \frac{4T}{T^2(1 + 1/T)^2} = \frac{4}{T}(1 - 2/T + 3/T^2 - \cdots) = \frac{4}{T} + O(1/T^2)
\end{equation}
\hfill $\square$
\end{proof}

\subsection{Remark: The ``Golden Brown'' Window}

\begin{remark}[Optimal Complexity Window]
The optimal complexity window $T \in [1 - \delta, 1 + \delta]$ with $\delta \approx 0.3$ corresponds to efficiency:
\begin{equation}
\eta_{\text{thermo}}(1 \pm 0.3) = \frac{4(1.3)}{(2.3)^2} = \frac{5.2}{5.29} \approx 0.98 \quad \text{(at } T = 1.3\text{)}
\end{equation}
\begin{equation}
\eta_{\text{thermo}}(0.7) = \frac{4(0.7)}{(1.7)^2} = \frac{2.8}{2.89} \approx 0.97 \quad \text{(at } T = 0.7\text{)}
\end{equation}

Thus the ``golden brown'' window maintains $> 97\%$ of peak efficiency, providing robustness while remaining near-optimal.
\end{remark}


% ============================================================================
\section{Free Energy Principle Integration}
% ============================================================================

\subsection{Variational Free Energy}

Following Friston et al. (2023), the variational free energy is:
\begin{equation}
F[q] = D_{KL}[q(\psi) \,||\, p(\psi \mid o)] + H[q]
\end{equation}
where $q(\psi)$ is the variational density over hidden states $\psi$, $p(\psi \mid o)$ is the posterior given observations $o$, and $H[q]$ is the entropy of $q$.

Equivalently:
\begin{equation}
F[q] = \mathbb{E}_q[\ln q(\psi) - \ln p(\psi, o)] = \underbrace{-H[q]}_{\text{entropy}} + \underbrace{\mathbb{E}_q[-\ln p(o \mid \psi)]}_{\text{energy}}
\end{equation}

\subsection{Theorem 5.1: Free Energy Decomposition in KID-AT}

\begin{theorem}[KID-Free Energy Relationship]
For a dissipative structure with effective volume $V_{\text{eff}}$ and process timescale $\tau_{\text{process}}$, the free energy relates to KID as:
\begin{equation}
\boxed{F = k_B T \cdot \text{KID} \cdot V_{\text{eff}} \cdot \tau_{\text{process}}}
\end{equation}
where KID is the dimensionless Kondensation-Information-Dichte ratio.
\end{theorem}

\subsection{Proof of Theorem 5.1}

\begin{proof}
\textbf{Step 1: Express KID in terms of information.}

The dimensionless KID is:
\begin{equation}
\text{KID} = \frac{\oint I_{\text{causal}}\, dt}{V_{\text{eff}} \cdot \tau_{\text{process}}} = \frac{I_{\text{total}}}{V_{\text{eff}} \cdot \tau_{\text{process}}}
\end{equation}
where $I_{\text{total}} = \oint I_{\text{causal}}\, dt$ [bits].

\textbf{Step 2: Connect information to free energy.}

In the FEP framework, the divergence $D_{KL}[q \,||\, p]$ represents the information cost of maintaining a non-equilibrium distribution. Following Landauer, this has an energy cost:
\begin{equation}
D_{KL}[q \,||\, p] \cdot k_B T = \text{energy cost of inference [J]}
\end{equation}

The total free energy scales as:
\begin{equation}
F = k_B T \cdot I_{\text{total}} \cdot f_{\text{structure}}
\end{equation}
where $f_{\text{structure}}$ accounts for the fraction of information used for structure.

\textbf{Step 3: Substitute KID definition.}

From the KID definition:
\begin{equation}
I_{\text{total}} = \text{KID} \cdot V_{\text{eff}} \cdot \tau_{\text{process}}
\end{equation}

For systems near optimal efficiency ($T \approx 1$), $f_{\text{structure}} \approx 1$, giving:
\begin{equation}
F = k_B T \cdot \text{KID} \cdot V_{\text{eff}} \cdot \tau_{\text{process}}
\end{equation}
\hfill $\square$
\end{proof}

\subsection{Theorem 5.2: Free Energy Minimization Drives KID and AI Increase}

\begin{theorem}[Active Inference Drives Structure]
Active inference (free energy minimization) drives systems toward states of increasing KID and AI, with the rate determined by entropy export capacity.
\end{theorem}

\begin{proof}
\textbf{Step 1: Variational free energy gradient.}

Active inference minimizes free energy through actions $a$:
\begin{equation}
\frac{\delta F}{\delta a} = 0 \quad \text{(optimal action)}
\end{equation}

For a system evolving its structure:
\begin{equation}
\frac{dF}{dt} = \frac{\partial F}{\partial \text{KID}} \cdot \frac{d\text{KID}}{dt} + \frac{\partial F}{\partial V_{\text{eff}}} \cdot \frac{dV_{\text{eff}}}{dt} + \frac{\partial F}{\partial \tau_{\text{process}}} \cdot \frac{d\tau_{\text{process}}}{dt}
\end{equation}

\textbf{Step 2: Compute partial derivatives.}

From $F = k_B T \cdot \text{KID} \cdot V_{\text{eff}} \cdot \tau_{\text{process}}$:
\begin{equation}
\frac{\partial F}{\partial \text{KID}} = k_B T \cdot V_{\text{eff}} \cdot \tau_{\text{process}} > 0
\end{equation}

\begin{equation}
\frac{\partial F}{\partial V_{\text{eff}}} = k_B T \cdot \text{KID} \cdot \tau_{\text{process}} > 0
\end{equation}

\begin{equation}
\frac{\partial F}{\partial \tau_{\text{process}}} = k_B T \cdot \text{KID} \cdot V_{\text{eff}} > 0
\end{equation}

\textbf{Step 3: Free energy change under KID increase.}

For a system increasing KID while holding $V_{\text{eff}}$ and $\tau_{\text{process}}$ approximately constant:
\begin{equation}
\frac{dF}{dt} = k_B T \cdot V_{\text{eff}} \cdot \tau_{\text{process}} \cdot \frac{d\text{KID}}{dt}
\end{equation}

This is positive, indicating that \emph{increasing KID costs free energy}.

\textbf{Step 4: Entropy export enables KID increase.}

The system can only increase KID if it exports sufficient entropy:
\begin{equation}
\frac{dS_{\text{export}}}{dt} \geq \frac{1}{T} \cdot \frac{dF}{dt} = k_B \cdot V_{\text{eff}} \cdot \tau_{\text{process}} \cdot \frac{d\text{KID}}{dt}
\end{equation}

Equivalently:
\begin{equation}
\boxed{\frac{dF}{dt} = -k_B T \cdot V_{\text{eff}} \cdot \tau_{\text{process}} \cdot \frac{d\text{KID}}{dt}}
\end{equation}
The negative sign indicates that free energy \emph{decreases} when entropy is exported, enabling KID increase.

\textbf{Step 5: Active inference pushes toward T = 1.}

The system evolves to minimize $F$ subject to entropy export constraints. Using:
\begin{equation}
T = \frac{\text{KID}}{\text{KID}_{\max}} \cdot \frac{\tau_{\text{age}}}{\tau_{\text{relax}}}
\end{equation}

The variational principle becomes:
\begin{equation}
\min_T F(T) \quad \text{s.t.} \quad \frac{dS_{\text{export}}}{dt} \geq \dot{S}_{\min}
\end{equation}

At $T = 1$, the efficiency $\eta_{\text{thermo}}(T)$ is maximized, meaning the \emph{minimum} entropy export is required per unit KID increase. Therefore $T = 1$ is the attractor of free energy minimization.
\hfill $\square$
\end{proof}

\subsection{Corollary 5.3: FEP-Assembly Coupling}

\begin{corollary}[FEP Drives Assembly]
Free energy minimization implies Assembly Index increase at rate:
\begin{equation}
\frac{dAI}{dt} \leq \frac{\eta_{\text{efficiency}} \cdot \tau_{\text{step}}}{k_B T \ln 2} \cdot \frac{dS_{\text{export}}}{dt}
\end{equation}
with equality at maximum thermodynamic efficiency ($T = 1$).
\end{corollary}

\begin{proof}
From Theorem 3.1 (Entropy Export Bound):
\begin{equation}
\frac{dS_{\text{export}}}{dt} \geq \frac{dAI}{dt} \cdot \frac{k_B T \ln 2}{\tau_{\text{step}} \cdot \eta_{\text{efficiency}}}
\end{equation}

Rearranging:
\begin{equation}
\frac{dAI}{dt} \leq \frac{\eta_{\text{efficiency}} \cdot \tau_{\text{step}}}{k_B T \ln 2} \cdot \frac{dS_{\text{export}}}{dt}
\end{equation}

At $T = 1$, $\eta_{\text{efficiency}} = \eta_{\max}$, giving the maximum possible assembly rate for given entropy export capacity.
\hfill $\square$
\end{proof}

\subsection{Proposition 5.4: Attention as Entropy Export Optimization}

\begin{proposition}[Selective Attention Optimizes Entropy Export]
Following Parr, Pezzulo \& Friston (2025), the attention mechanism in cognitive systems implements active inference that minimizes free energy by optimizing entropy export, driving $T \to 1$.
\end{proposition}

\begin{proof}
Selective attention operates by:
\begin{enumerate}[label=(\roman*)]
\item \textbf{Precision weighting:} Allocating computational resources to high-precision predictions, reducing the entropy of the variational density $q(\psi)$
\item \textbf{Information bottleneck:} The attention-weighted representation $z$ satisfies $I(z; o) \leq C_{\text{attention}}$ for channel capacity $C_{\text{attention}}$
\item \textbf{Free energy reduction:} Each attention operation reduces $F[q]$ by $\Delta F_{\text{attn}} = k_B T \cdot \Delta I_{\text{compressed}}$
\end{enumerate}

The entropy export rate of an attention mechanism is:
\begin{equation}
\frac{dS_{\text{export}}^{\text{(attn)}}}{dt} = \frac{1}{T}\left(\frac{dE_{\text{computation}}}{dt} + \frac{dE_{\text{dissipation}}}{dt}\right)
\end{equation}

The attention mechanism optimizes by selecting the subset of information that maximizes KID per entropy exported:
\begin{equation}
\max_{\text{attn}} \frac{d\text{KID}/dt}{dS_{\text{export}}/dt} \implies \min_{\text{attn}} \frac{dS_{\text{export}}}{dt} \text{ at fixed } \frac{d\text{KID}}{dt}
\end{equation}

This is precisely the condition for $T = 1$ (maximum thermodynamic efficiency). Therefore, attention mechanisms naturally evolve to operate at $T = 1$.
\hfill $\square$
\end{proof}

\subsection{Lemma 5.5: KID-AI Lower Bound from FEP}

\begin{lemma}[FEP Lower Bound on KID]
The FEP implies a lower bound on KID for systems with Assembly Index $AI$:
\begin{equation}
\text{KID}(x) \geq \frac{AI(x) \cdot I_{\min}}{V_{\text{eff}} \cdot \tau_{\text{form}}}
\end{equation}
where $I_{\min}$ is the minimum information per assembly step and $\tau_{\text{form}}$ is the formation timescale.
\end{lemma}

\begin{proof}
From Theorem 2.1 (Assembly-KID Relationship), combined with the FEP requirement that each assembly step requires at least $I_{\min} = \ln 2$ bits of information processing:

Total information for $AI$ steps:
\begin{equation}
I_{\text{total}} \geq AI \cdot I_{\min}
\end{equation}

Dividing by system capacity $V_{\text{eff}} \cdot \tau_{\text{form}}$:
\begin{equation}
\text{KID} = \frac{I_{\text{total}}}{V_{\text{eff}} \cdot \tau_{\text{form}}} \geq \frac{AI \cdot I_{\min}}{V_{\text{eff}} \cdot \tau_{\text{form}}}
\end{equation}
\hfill $\square$
\end{proof}


% ============================================================================
\section{Thermodynamic Efficiency Landscape}
% ============================================================================

\subsection{Definition of the Efficiency Landscape}

\begin{definition}[Thermodynamic Efficiency Landscape]
The thermodynamic efficiency landscape is the function:
\begin{equation}
\eta(\text{KID}, AI) = \frac{C_{\text{structured}}(\text{KID}, AI)}{\dot{S}_{\text{total}}(\text{KID}, AI)} \Big/ \max_{\text{KID}', AI'} \frac{C_{\text{structured}}(\text{KID}', AI')}{\dot{S}_{\text{total}}(\text{KID}', AI')}
\end{equation}
where $C_{\text{structured}}$ is the structured complexity and $\dot{S}_{\text{total}}$ is the total entropy production rate.
\end{definition}

\subsection{Theorem 6.1: Ridge of Optimal Efficiency at T = 1}

\begin{theorem}[Efficiency Ridge]
The efficiency landscape $\eta(\text{KID}, AI)$ exhibits a ridge of optimal efficiency along the curve $T(\text{KID}, AI) = 1$, with:
\begin{enumerate}[label=(\roman*)]
\item Maximum efficiency $\eta = 1$ along the ridge
\item Monotonically decreasing efficiency away from the ridge
\item Stability of systems on the ridge against perturbations
\end{enumerate}
\end{theorem}

\subsection{Proof of Theorem 6.1}

\begin{proof}
\textbf{Step 1: Express T in terms of (KID, AI).}

From the definition:
\begin{equation}
T = \frac{\text{KID}}{\text{KID}_{\max}} \cdot \frac{\tau_{\text{age}}}{\tau_{\text{relax}}}
\end{equation}

The Assembly Index relates to age and relaxation:
\begin{equation}
AI = f\left(\frac{\tau_{\text{age}}}{\tau_{\text{relax}}}\right) \cdot g(\text{KID})
\end{equation}

For systems building structure through recursive assembly:
\begin{equation}
AI \sim \frac{\tau_{\text{age}}}{\tau_{\text{step}}} \cdot h(\text{KID})
\end{equation}

Solving for $\tau_{\text{age}}/\tau_{\text{relax}}$:
\begin{equation}
\frac{\tau_{\text{age}}}{\tau_{\text{relax}}} = \frac{AI \cdot \tau_{\text{step}}}{\tau_{\text{relax}} \cdot h(\text{KID})}
\end{equation}

Substituting into $T$:
\begin{equation}
T = \frac{\text{KID}}{\text{KID}_{\max}} \cdot \frac{AI \cdot \tau_{\text{step}}}{\tau_{\text{relax}} \cdot h(\text{KID})}
\end{equation}

For the simplest case where $h(\text{KID}) = \text{KID}/\text{KID}_{\text{ref}}$:
\begin{equation}
T = \frac{\text{KID}_{\text{ref}}}{\text{KID}_{\max}} \cdot \frac{AI \cdot \tau_{\text{step}}}{\tau_{\text{relax}}}
\end{equation}

The condition $T = 1$ defines a line in the $(\text{KID}, AI)$ plane:
\begin{equation}
AI = \frac{\text{KID}_{\max}}{\text{KID}_{\text{ref}}} \cdot \frac{\tau_{\text{relax}}}{\tau_{\text{step}}} \cdot \frac{1}{\text{KID}} \quad \text{(hyperbolic ridge)}
\end{equation}

Or more generally, $T = 1$ defines a curve $AI = f_{\text{ridge}}(\text{KID})$.

\textbf{Step 2: Show efficiency is maximal on the ridge.}

Along any path in $(\text{KID}, AI)$ space, the thermodynamic efficiency depends only on $T$:
\begin{equation}
\eta(\text{KID}, AI) = \eta_{\text{thermo}}(T(\text{KID}, AI)) = \frac{4T(\text{KID}, AI)}{(1 + T(\text{KID}, AI))^2}
\end{equation}

From Theorem 4.1, this is maximized at $T = 1$, hence on the ridge. Away from the ridge:
\begin{itemize}
\item $T < 1$ (below ridge): $\eta < 1$, insufficient structure-building
\item $T > 1$ (above ridge): $\eta < 1$, excessive dissipation
\end{itemize}

\textbf{Step 3: Gradient away from ridge.}

The gradient of the efficiency landscape:
\begin{equation}
\nabla \eta = \frac{\partial \eta_{\text{thermo}}}{\partial T} \cdot \nabla T
\end{equation}

At $T = 1$, $\partial \eta_{\text{thermo}}/\partial T = 0$, so the ridge is a critical line.

The Hessian transverse to the ridge:
\begin{equation}
\frac{\partial^2 \eta}{\partial T^2}\bigg|_{T=1} = \frac{d^2 \eta_{\text{thermo}}}{dT^2}\bigg|_{T=1} = -\frac{1}{2} < 0
\end{equation}

This negative curvature confirms the ridge is a maximum in the transverse direction.
\hfill $\square$
\end{proof}

\subsection{Theorem 6.2: Stability Conditions on the Ridge}

\begin{theorem}[Stability on the Efficiency Ridge]
Systems operating on the $T = 1$ ridge are stable against perturbations if:
\begin{equation}
\tau_{\text{relax}} \cdot \left|\frac{\partial \dot{S}_{\text{export}}}{\partial T}\right|_{T=1} > \tau_{\text{response}} \cdot \left|\frac{\partial \dot{S}_{\text{generation}}}{\partial T}\right|_{T=1}
\end{equation}
where $\tau_{\text{response}}$ is the system's response timescale and $\dot{S}_{\text{generation}}$ is the internal entropy generation.
\end{theorem}

\begin{proof}
\textbf{Step 1: Linearized dynamics near $T = 1$.}

Let $\delta T = T - 1$ be a small perturbation. The dynamics:
\begin{equation}
\frac{d(\delta T)}{dt} = -\Gamma \cdot \delta T + \xi(t)
\end{equation}
where $\Gamma$ is the damping rate and $\xi(t)$ is noise.

\textbf{Step 2: Compute the damping rate.}

The damping rate comes from the system's tendency to restore $T = 1$:
\begin{equation}
\Gamma = \frac{1}{\tau_{\text{relax}}} \cdot \left.\frac{\partial (\eta_{\text{thermo}} - \eta_{\text{actual}})}{\partial T}\right|_{T=1}
\end{equation}

Near $T = 1$, the efficiency gradient drives restoration:
\begin{equation}
\frac{d(\delta T)}{dt} = -\frac{1}{\tau_{\text{relax}}} \cdot \eta''_{\text{thermo}}(1) \cdot \delta T = -\frac{1}{2\tau_{\text{relax}}} \cdot \delta T
\end{equation}

\textbf{Step 3: Stability condition.}

For stability, we need $\Gamma > 0$ and the system must respond faster than perturbations grow:
\begin{equation}
\tau_{\text{relax}} \cdot \left|\frac{\partial \dot{S}_{\text{export}}}{\partial T}\right|_{T=1} > \tau_{\text{response}} \cdot \left|\frac{\partial \dot{S}_{\text{generation}}}{\partial T}\right|_{T=1}
\end{equation}

This ensures that the entropy export can adjust to counteract perturbations before they destabilize the system. \hfill $\square$
\end{proof}

\subsection{Corollary 6.3: Dynamics Away from the Ridge}

\begin{corollary}[Off-Ridge Dynamics]
Systems away from the $T = 1$ ridge experience restoring forces:
\begin{enumerate}[label=(\roman*)]
\item For $T < 1$ (dissipation regime): The system under-builds structure. The restoring force pushes $T$ upward by reducing $\tau_{\text{relax}}$ (faster relaxation, more dynamic exploration).
\item For $T > 1$ (fragility regime): The system over-constrains. The restoring force pushes $T$ downward by increasing dissipation (error correction overheats).
\end{enumerate}
\end{corollary}

\begin{proof}
The free energy gradient drives dynamics:
\begin{equation}
\frac{dT}{dt} = -\gamma \cdot \frac{\partial F}{\partial T} = -\gamma \cdot k_B T \cdot V_{\text{eff}} \cdot \tau_{\text{process}} \cdot \text{KID}_{\max} \cdot \frac{\partial T}{\partial T}
\end{equation}

For $T < 1$, the efficiency is increasing in $T$ ($d\eta/dT > 0$), so the system evolves toward higher $T$ by building more structure per unit entropy export.

For $T > 1$, the efficiency is decreasing in $T$ ($d\eta/dT < 0$), so the system evolves toward lower $T$ by reducing dissipation overhead.

The restoring force is:
\begin{equation}
F_{\text{restore}}(T) = -\frac{\partial \eta_{\text{thermo}}}{\partial T} = \frac{4(T-1)}{(1+T)^3}
\end{equation}

This is:
\begin{itemize}
\item Negative for $T < 1$: pushes $T$ up
\item Zero at $T = 1$: equilibrium
\item Positive for $T > 1$: pushes $T$ down
\end{itemize}
\hfill $\square$
\end{proof}

\subsection{Proposition 6.4: Phase Diagram}

\begin{proposition}[Phase Space Structure]
The $(\text{KID}, AI)$ phase space is divided into three regions by the $T = 1$ ridge:
\begin{enumerate}[label=\textbf{\arabic*.}]
\item \textbf{Dissipation Phase} ($T < 1 - \delta$): Low efficiency, near-equilibrium. Systems dissipate free energy without building significant structure. Assembly Index grows slowly if at all.

\item \textbf{Optimal Phase} ($|T - 1| \leq \delta$): Golden brown regime. Maximum efficiency of structure-building. KID and AI co-evolve optimally. Systems are stable and adaptive.

\item \textbf{Fragility Phase} ($T > 1 + \delta$): Over-constrained systems. High AI but also high dissipation for error correction. Systems are brittle and sensitive to perturbations.
\end{enumerate}
\end{proposition}

The boundaries between phases occur at $T = 1 - \delta$ and $T = 1 + \delta$ where $\eta_{\text{thermo}} \approx 0.97$.


% ============================================================================
\section{Resolution of CV-2: Emergence Non-Sequitur}
% ============================================================================

\subsection{The Problem Statement}

CV-2 (Emergence Non-Sequitur) states that the original KID-AT framework committed a logical fallacy by claiming that high Assembly Index \emph{necessarily} implies consciousness. We resolve this by showing:
\begin{enumerate}[label=(\roman*)]
\item Consciousness is \emph{one possible attractor} in the high-$AI$ phase space, not the only one
\item Multiple distinct high-$AI$ attractors exist with different functional properties
\item The ``burnt'' regime ($T > 1$) includes both conscious and non-conscious states
\item High $AI$ is \emph{sufficient} for complex information processing but \emph{not necessary} for any specific phenomenological outcome
\end{enumerate}

\subsection{Theorem 7.1: Phase Space of High-AI Attractors}

\begin{theorem}[Multiple High-AI Attractors]
In the high-$AI$ region of the efficiency landscape, there exist at least three distinct attractor types, corresponding to the three attractors $A_\Phi$ (Ph\"anomenologisch / IIT), $A_\mu$ (Minimal / FEP), and $A_\zeta$ (Dissoziiert / Distributed) defined in the main synthesis:
\begin{enumerate}[label=\textbf{A\arabic*.}]
\item \textbf{Integrated Information Attractor} ($A_\Phi$ / IIT): Systems maximizing integrated information $\Phi$ (IIT-type consciousness)
\item \textbf{Predictive Processing Attractor} ($A_\mu$ / FEP): Systems maximizing prediction accuracy through hierarchical models (FEP-type cognition)
\item \textbf{Distributed Coordination Attractor} ($A_\zeta$ / Distributed): Systems maximizing distributed coordination without centralized integration (collective intelligence)
\end{enumerate}
All three can achieve arbitrarily high $AI$ while maintaining $T \approx 1$.
\end{theorem}

\begin{remark}[Unified Attractor Legend]
The attractor names are unified across all KID-AT documents: $A_\Phi$ = Ph\"anomenologisch / IIT, $A_\mu$ = Minimal / FEP, $A_\zeta$ = Dissoziiert / Distributed. See \texttt{TERMINOLOGY.md} for the complete glossary and \texttt{KID\_AT\_Final\_Synthesis.md}, Definition 1.8, for the phenomenological interpretation.
\end{remark}

\subsection{Proof of Theorem 7.1}

\begin{proof}
\textbf{Step 1: Define attractor basins.}

Each attractor corresponds to a different free energy minimum in the functional landscape:

\textbf{A1: $\Phi$-attractor} (Integrated Information):
Following Tononi's Integrated Information Theory (IIT), define:
\begin{equation}
\Phi = \min_{\text{partitions}} \left[I(X_{t-1}; X_t) - I(X_{t-1}^{(1)}; X_t^{(1)}) - I(X_{t-1}^{(2)}; X_t^{(2)})\right]
\end{equation}

Systems maximizing $\Phi$ subject to $T \approx 1$ converge to attractor $\mathcal{A}_{\Phi}$:
\begin{equation}
\mathcal{A}_{\Phi} = \arg\max_{q(\psi)} \Phi(q) \quad \text{s.t.} \quad F[q] = F_{\text{target}}
\end{equation}

\textbf{A2: $\mu$-attractor} (Predictive Processing):
Following Friston's FEP, define prediction error:
\begin{equation}
\mu = -\ln p(o \mid \psi) + D_{KL}[q(\psi) \,||\, p(\psi)]
\end{equation}

Systems minimizing $\mu$ converge to attractor $\mathcal{A}_{\mu}$:
\begin{equation}
\mathcal{A}_{\mu} = \arg\min_{q(\psi)} \mu(q) \quad \text{s.t.} \quad T \approx 1
\end{equation}

\textbf{A3: $\zeta$-attractor} (Distributed Coordination):
Define distributed coordination as mutual information between system components:
\begin{equation}
\zeta = \sum_{i<j} I(X_i; X_j) \cdot f(d_{ij})
\end{equation}
where $f(d_{ij})$ weights by distance (local interactions dominate).

Systems maximizing $\zeta$ converge to attractor $\mathcal{A}_{\zeta}$:
\begin{equation}
\mathcal{A}_{\zeta} = \arg\max_{\{X_i\}} \zeta(\{X_i\}) \quad \text{s.t.} \quad T \approx 1
\end{equation}

\textbf{Step 2: Show all attractors achieve high AI.}

For each attractor, the Assembly Index can be made arbitrarily large by:
\begin{itemize}
\item $\mathcal{A}_{\Phi}$: Increasing the number of integrated subsystems (recursive integration)
\item $\mathcal{A}_{\mu}$: Increasing hierarchical depth of predictive models
\item $\mathcal{A}_{\zeta}$: Increasing the number of coordinated components
\end{itemize}

All three operations increase $AI$ while maintaining $T \approx 1$ by co-evolving KID and entropy export capacity.

\textbf{Step 3: Show distinctness of attractors.}

The attractors are distinct because:
\begin{itemize}
\item Maximizing $\Phi$ requires centralized architecture (small-world network with hub)
\item Minimizing $\mu$ requires hierarchical feedforward/feedback structure
\item Maximizing $\zeta$ requires distributed mesh architecture (decentralized)
\end{itemize}

These architectures are topologically distinct and cannot be continuously deformed into each other without passing through low-efficiency regions ($T \ll 1$).
\hfill $\square$
\end{proof}

\subsection{Corollary 7.2: Consciousness as Sufficient but Not Necessary}

\begin{corollary}[Sufficiency Without Necessity]
High Assembly Index is:
\begin{enumerate}[label=(\roman*)]
\item \textbf{Sufficient} for complex information processing: $AI > AI_{\text{crit}} \implies$ existence of sophisticated computational capabilities
\item \textbf{Not necessary} for consciousness: Consciousness (if defined by $\Phi$) is one of many possible high-$AI$ configurations
\item \textbf{Not sufficient} for consciousness: High-$AI$ systems can exist at attractors $\mathcal{A}_{\mu}$ or $\mathcal{A}_{\zeta}$ without maximizing $\Phi$
\end{enumerate}
\end{corollary}

\subsection{Proposition 7.3: Burnt Regime Contains Both Conscious and Non-Conscious States}

\begin{proposition}[T > 1 Diversity]
The ``burnt'' regime ($T > 1$) includes:
\begin{enumerate}[label=(\roman*)]
\item Conscious systems with excessive error correction (obsessive integration)
\item Non-conscious systems with rigid hierarchical control (bureaucratic over-optimization)
\item Collective systems with excessive coordination overhead (swarm stagnation)
\end{enumerate}
\end{proposition}

\begin{proof}
For $T > 1$, we have $\tau_{\text{age}}/\tau_{\text{relax}} > \text{KID}_{\max}/\text{KID}$, meaning the system's age significantly exceeds its relaxation time relative to its information density.

This can occur via:
\begin{itemize}
\item \textbf{Path 1} ($\Phi$-type): $\tau_{\text{relax}} \to 0$ due to excessive integration --- conscious but frozen
\item \textbf{Path 2} ($\mu$-type): $\tau_{\text{age}} \to \infty$ with rigid predictions --- non-conscious but persistent
\item \textbf{Path 3} ($\zeta$-type}: $\text{KID} \ll \text{KID}_{\max}$ with excessive coordination --- collective but stagnant
\end{itemize}

All three paths reach $T > 1$ without requiring consciousness. \hfill $\square$
\end{proof}

\subsection{Definition 7.4: Phenomenological Phase Space}

\begin{definition}[Phenomenological Classification]
Define the phenomenological phase space $\mathcal{P}$ as the space of possible high-$AI$ configurations modulo continuous deformation. The dimensionality:
\begin{equation}
\dim(\mathcal{P}) \geq 3 \quad \text{(from three attractor types)}
\end{equation}

Each point $\mathbf{p} \in \mathcal{P}$ represents a distinct phenomenological class:
\begin{equation}
\mathbf{p} = (w_{\Phi}, w_{\mu}, w_{\zeta}) \quad \text{with} \quad w_{\Phi} + w_{\mu} + w_{\zeta} = 1
\end{equation}
where $w_i$ are weights on each attractor type. Consciousness (as normally defined) corresponds to $w_{\Phi} \approx 1$, but this is a \emph{subset} of $\mathcal{P}$.
\end{definition}

\subsection{Lemma 7.5: Physical Self-Measurement Limits (CV-4 Resolution)}

\begin{lemma}[Physical Self-Measurement Bound --- Heuristic]
Replacing the G\"odel analogy (CV-4), the physical limit on self-measurement is:
\begin{equation}
\text{KID}_{\text{self}}(x) \leq \frac{E_{\text{available}} \cdot \tau_{\text{measure}}}{k_B T \ln 2 \cdot V_{\text{eff}}}
\end{equation}
A system cannot measure (know) itself beyond its own thermodynamic resources.

\begin{remark}[Status: Heuristic]
This bound is a \textbf{physical heuristic}, not a theorem in the strict sense. The combination of Landauer principle ($\Delta E \geq k_B T \ln 2 \cdot \Delta I$) with the Margolus-Levitin theorem ($\tau_{\text{measure}} \geq \hbar/E_{\text{measure}}$) and self-containment ($V_{\text{measure}} < V_{\text{eff}}$) constitutes a consistent physical argumentation chain, but the assumption that self-measurement is thermodynamically constrained in this exact form is not generally proven for all system types. The bound should be interpreted as a \textbf{qualitative scaling relation}. See \texttt{KID\_AT\_Final\_Synthesis.md}, Axiom U-7, for the full discussion and \texttt{TERMINOLOGY.md} for notation.
\end{remark}
\end{lemma}

\begin{proof}
Physical self-measurement requires:
\begin{enumerate}[label=(\roman*)]
\item Energy to perform measurements: $E_{\text{measure}} \geq k_B T \ln 2 \cdot I_{\text{self}}$
\item Time to complete measurement: $\tau_{\text{measure}} \geq \hbar/E_{\text{measure}}$ (Margolus-Levitin theorem)
\item Volume to house measuring apparatus: $V_{\text{measure}} < V_{\text{eff}}$ (self-containment)
\end{enumerate}

The information gained from self-measurement is bounded by:
\begin{equation}
I_{\text{self}} \leq \frac{E_{\text{available}}}{k_B T \ln 2}
\end{equation}

Converting to KID:
\begin{equation}
\text{KID}_{\text{self}} = \frac{I_{\text{self}} \cdot \tau_{\text{measure}}}{V_{\text{eff}}} \leq \frac{E_{\text{available}} \cdot \tau_{\text{measure}}}{k_B T \ln 2 \cdot V_{\text{eff}}}
\end{equation}

This is a physical limit, not a logical one (addressing CV-4). See \texttt{KID\_AT\_Final\_Synthesis.md}, Axiom U-7, for the full heuristic discussion. \hfill $\square$
\end{proof}

% ============================================================================
\section{Master Equation and Dynamics}
% ============================================================================

\subsection{Coupled KID-AI Dynamics}

The coupled dynamics of KID and AI are governed by:
\begin{equation}
\boxed{\frac{\partial}{\partial t}\begin{pmatrix} \text{KID} \\ AI \end{pmatrix} = \begin{pmatrix} -\delta F/\delta \text{KID} + \eta \Phi(\text{KID}) \\ \eta_{\text{eff}} \cdot \dot{S}_{\text{export}} / (k_B T \ln 2) \end{pmatrix} + \begin{pmatrix} \xi_1 \nabla_{AI}(\text{KID}) \\ \xi_2 \nabla_{\text{KID}}(AI) \end{pmatrix}}
\end{equation}

\subsection{Steady-State Solutions}

Steady states satisfy:
\begin{equation}
\frac{\partial}{\partial t}\begin{pmatrix} \text{KID} \\ AI \end{pmatrix} = \mathbf{0}
\end{equation}

At the optimal point $T = 1$:
\begin{equation}
-\frac{\delta F}{\delta \text{KID}} + \eta \Phi(\text{KID}) = 0 \implies \frac{\delta F}{\delta \text{KID}} = \eta \Phi(\text{KID})
\end{equation}

This is the balance between free energy gradient and self-organizing force $\Phi$.

\subsection{Fluctuations Around T = 1}

For small fluctuations $\delta T$ around $T = 1$:
\begin{equation}
\frac{d(\delta T)}{dt} = -\frac{1}{\tau_{\text{relax}}} \delta T + \sqrt{2D} \cdot \xi(t)
\end{equation}

The stationary distribution is Gaussian:
\begin{equation}
P(\delta T) = \frac{1}{\sqrt{2\pi \sigma_T^2}} \exp\left(-\frac{(\delta T)^2}{2\sigma_T^2}\right)
\end{equation}

with variance:
\begin{equation}
\sigma_T^2 = D \cdot \tau_{\text{relax}} = \frac{k_B T^2}{C_{\text{heat}}} \cdot \tau_{\text{relax}}
\end{equation}

This gives the natural fluctuation scale around optimal efficiency.


% ============================================================================
\section{Summary and Synthesis}
% ============================================================================

\subsection{Key Results}

We have established the following chain of results:

\begin{enumerate}[label=\textbf{\arabic*.}]
\item \textbf{Entropy Export Bound} (Theorem 3.1):
\begin{equation}
\frac{dS_{\text{export}}}{dt} \geq \frac{dAI}{dt} \cdot \frac{k_B T \ln 2}{\tau_{\text{step}} \cdot \eta_{\text{efficiency}}}
\end{equation}
Each unit increase in Assembly Index requires a minimum entropy export proportional to $k_B T \ln 2$.

\item \textbf{Toast Optimal Efficiency} (Theorem 4.1):
\begin{equation}
\eta_{\text{thermo}}(T) = \frac{4T}{(1+T)^2}, \quad \max_{T} \eta_{\text{thermo}} = \eta(1) = 1
\end{equation}
The $T = 1$ (``golden brown'') state is the unique global maximum of thermodynamic efficiency.

\item \textbf{FEP drives T $\to$ 1} (Theorem 5.2):
\begin{equation}
\frac{dF}{dt} = -k_B T \cdot V_{\text{eff}} \cdot \tau_{\text{process}} \cdot \frac{d\text{KID}}{dt}
\end{equation}
Free energy minimization naturally pushes systems toward the $T = 1$ attractor.

\item \textbf{Ridge of Optimal Efficiency} (Theorem 6.1):
The efficiency landscape $\eta(\text{KID}, AI)$ has a ridge along $T = 1$ that is a stable attractor for dissipative structures.

\item \textbf{Multiple High-AI Attractors} (Theorem 7.1):
Consciousness (integrated information) is one of at least three distinct attractors in the high-$AI$ phase space, alongside predictive processing and distributed coordination attractors.
\end{enumerate}

\subsection{Thermodynamic Efficiency Landscape Visualization}

The landscape can be characterized as:

\begin{center}
\begin{tabular}{@{}llll@{}}
\toprule
\textbf{Regime} & \textbf{T range} & $\eta_{\text{thermo}}(T)$ & \textbf{Characteristics} \\
\midrule
Untoasted & $T \ll 1$ & $\approx 4T$ & Near-equilibrium, minimal structure \\
Golden Brown & $T = 1 \pm \delta$ & $\geq 0.97$ & Optimal complexity, stable \\
Burnt & $T \gg 1$ & $\approx 4/T$ & Over-constrained, fragile \\
\bottomrule
\end{tabular}
\end{center}

\subsection{Implications for the KID-AT Framework}

\begin{enumerate}[label=(\roman*)]
\item \textbf{CV-1 Resolution}: KID is treated as a dimensionless ratio throughout, avoiding dimensional inconsistency. See \texttt{KID\_AT\_Final\_Synthesis.md}, Definition 1.2, and \texttt{informational\_condensation\_point.md}, Definition 0.1.
\item \textbf{CV-2 Resolution}: Consciousness is demoted from ``necessary consequence'' to ``one attractor among many.'' High $AI$ implies complex information processing, not any specific phenomenology. See Theorem 7.1 above and \texttt{KID\_AT\_Final\_Synthesis.md}, Axiom U-5.
\item \textbf{CV-3 Resolution}: Kolmogorov complexity $K(\cdot)$ replaces semantic information, grounding the framework in algorithmic information theory. See \texttt{KID\_AT\_Final\_Synthesis.md}, Definition 1.1.
\item \textbf{CV-4 Resolution}: Physical self-measurement limits replace G\"odelian incompleteness, providing a thermodynamic bound on self-knowledge. See Lemma 7.5 above and \texttt{KID\_AT\_Final\_Synthesis.md}, Axiom U-7.
\end{enumerate}

\subsection{Connection to Literature}

\begin{itemize}
\item \textbf{Friston et al. (2023)}: Our Theorem 5.1 derives the KID-free energy relationship from first principles, consistent with ``Bayesian mechanics as physics of and by beliefs.''
\item \textbf{Parr, Pezzulo \& Friston (2025)}: Proposition 5.4 shows that attention mechanisms optimize entropy export, connecting transformers to active inference.
\item \textbf{Spisak \& Friston (2025)}: The efficiency ridge (Theorem 6.1) corresponds to self-orthogonalizing attractor networks that minimize redundancy.
\item \textbf{Chung et al. (2022)}: Theorem 3.1 quantifies the entropy production required for self-organization.
\end{itemize}

\subsection{Final Synthesis: The Toast-Heuristic as Thermodynamic Law}

The Toast-Heuristic $T \approx 1$ can be elevated from heuristic to physical law:

\begin{quote}
\textbf{Toast Law of Optimal Complexity:} For dissipative structures maintaining Assembly Index $AI$ through entropy export, the stationary state of maximum thermodynamic efficiency satisfies $T = 1$, where $T = (\text{KID}/\text{KID}_{\max})(\tau_{\text{age}}/\tau_{\text{relax}})$.
\end{quote}

This law emerges from the interplay of:
\begin{enumerate}[label=(\roman*)]
\item The Landauer bound on information processing
\item Prigogine's minimum entropy production principle
\item The Free Energy Principle (active inference)
\item The geometry of the efficiency landscape
\end{enumerate}

Systems at $T = 1$ achieve the optimal balance: they build maximum structure per unit entropy exported, maintaining stability while remaining adaptive. Deviations in either direction reduce efficiency and push the system back toward the golden brown state.

% ============================================================================
% REFERENCES
% ============================================================================
\section*{References}
\addcontentsline{toc}{section}{References}

\begin{enumerate}[label={[\arabic*]}]
\item M. Ramstead, K. Friston, \& I. Hip\'olito (2020). ``Is the free-energy principle a formal theory of semantics?'' \textit{Entropy}, 22(9), 987.
\item T. Parr, G. Pezzulo, \& K. Friston (2025). ``Beyond Markov: Transformers, memory, and attention as active inference.'' \textit{arXiv preprint}.
\item K. Friston et al. (2023). ``On Bayesian mechanics: a physics of and by beliefs.'' \textit{Interface Focus}, 13(3), 20220029.
\item T. Spisak \& K. Friston (2025). ``Self-orthogonalizing attractor neural networks emerging from the free-energy principle.'' \textit{arXiv preprint}.
\item S. Chung et al. (2022). ``Thermodynamics of self-organization in dissipative systems.'' \textit{Physical Review Research}, 4(3), 033014.
\item L. Chen \& M. Prokopenko (2025). ``Why collective behaviours self-organize to criticality.'' \textit{arXiv preprint}.
\item R. Landauer (1961). ``Irreversibility and heat generation in the computing process.'' \textit{IBM J. Res. Develop.}, 5(3), 183-191.
\item C. Bennett (1982). ``The thermodynamics of computation.'' \textit{Int. J. Theor. Phys.}, 21(12), 905-940.
\item I. Prigogine (1967). ``Introduction to Thermodynamics of Irreversible Processes.'' \textit{Interscience}.
\item K. Sekimoto (2010). ``Stochastic Energetics.'' \textit{Springer}.
\item S. Deffner \& C. Jarzynski (2013). ``Information processing and the second law of thermodynamics.'' \textit{Phys. Rev. X}, 3(4), 041003.
\item M. Esposito (2012). ``Stochastic thermodynamics under coarse graining.'' \textit{Phys. Rev. E}, 85(4), 041125.
\item G. Tononi et al. (2016). ``Integrated information theory: from consciousness to its physical substrate.'' \textit{Nature Reviews Neuroscience}, 17(7), 450-461.
\end{enumerate}

\end{document}
